\documentclass[10pt,a4paper]{amsart}
\usepackage[margin=19mm]{geometry}
\usepackage{amsmath,amssymb,mathtools}
\usepackage[scr=rsfs,cal=euler]{mathalfa}
\usepackage{xfrac}
\usepackage[unicode,hidelinks,bookmarks=false]{hyperref}
\newcommand{\ie}{i.e.\ }
\newcommand{\cf}{cf.\ }
\newcommand{\resp}{respectively\ }
\newcommand{\Subsection}{\subsection}
\providecommand{\nobreakdash}{\nobreak\hbox{-}}
\setlength{\emergencystretch}{2em}
\multlinegap0pt
\usepackage{mathtools}
\usepackage{bm}
\usepackage[shortlabels]{enumitem}
\usepackage{xcolor}
\usepackage{tikz}
\newtheorem{theorem}{Theorem}
\newtheorem{lemma}{Lemma}
\usepackage{cleveref}
\crefname{theorem}{Theorem}{Theorems}
\crefname{lemma}{Lemma}{Lemmas}
\newcommand{\Dplus}{D^+}
\newcommand{\R}{\mathbb R}
\newcommand{\T}{\mathbb T}
\newcommand{\dist}{\operatorname{dist}}
\newcommand{\e}{\varepsilon}
\title{The Physical Cutoff Does Not Restore Homogenization: Phase-Dependent Burning in the Strain G-Equation}
\author{Michele Caprio}
\date{Source: arXiv:2608.15337v1 (CC BY 4.0)}
\begin{document}
\maketitle

\noindent\emph{Source and licence.} The Physical Cutoff Does Not Restore Homogenization: Phase-Dependent Burning in the Strain G-Equation. Version arXiv:2608.15337v1 (CC BY 4.0), distributed by arXiv under the Creative Commons Attribution 4.0 International licence, http://creativecommons.org/licenses/by/4.0/, as recorded in the arXiv licence metadata for this exact version and shown on the abstract page https://arxiv.org/abs/2608.15337v1. Full licence notice: \texttt{licenses/arxiv-2608.15337-CC-BY-4.0.txt}.\par\smallskip

\noindent\textit{Editorial scope.} Counted unit: the article's theorem barrier (source0.tex lines 880--911). The article's complete original proof of it is delivered with it (source0.tex lines 2315--2449, one \texttt{proof} environment of 135 lines, which the article places away from the statement). No mathematics is written, repaired or re-derived: the statement and the proof are the article's own lines, in the article's own notation, and no retained line is substituted or re-flowed.  Retained uncounted as the source-local context the proof needs: lines 2888--2888; lines 2905--2918.  Nothing was omitted from any retained span, so the package carries no omission record.  No retained line contains a citation, so the wrapper carries no bibliography and no bibliography entry.  No retained line contains a figure environment, an image-inclusion command or drawing code, so no image is delivered and the package declares no asset dependency.  Editorial formatting only, in addition: an AMS \texttt{amsart} preamble that restates the article's own declarations and macros used by the retained lines, namely the environments \texttt{theorem}, \texttt{lemma} on separate counters, together with the standard packages those lines need.  The retained environments print in this document as Theorem 1, Lemma 1 in order of appearance, whereas the article prints the same items with the numbering of its own full document.  The complete native source of the article as distributed in the arXiv e-print for this exact version is bundled unchanged as \texttt{sources/207735/source0.tex}.
\subsection{Proximal Normals, Distance Derivatives, and Incompressibility}\label{app:proximal-background}

\begin{lemma}[Upper Derivative of the Distance]\label{lem:distance-Dini-background}
Let $y(\cdot)$ be differentiable at $t$, let $y(t)\notin D$, and suppose
\[
  \rho(t)\coloneqq \dist(y(t),D)<\frac L2.
\]
Choose a nearest point $x\in D$ and write
\[
  y(t)=x+\rho(t)\nu,
  \qquad |\nu|=1.
\]
Then
  $D^+\rho(t)\le \dot y(t)\cdot\nu$.
The same inequality holds almost everywhere for absolutely continuous curves.
\end{lemma}

\begin{theorem}[Rigidity of Outward Barrier Certificates]
\label{thm:no barrier}
Let $n\ge2$ and $L>0$, and write
$\T_L^n\coloneqq\mathbb R^n/L\mathbb Z^n$.  Let
$V\in C^2(\T_L^n;\mathbb R^n)$ be divergence free, let $d>0$, and
define
\[
  H(x,q)
  \coloneqq
  \left(
    1+d\frac{q^{\mathsf T}S_V(x)q}{|q|^2}
  \right)_+|q|
  +V(x)\cdot q,
  \qquad H(x,0)\coloneqq0.
\]
We denote by $\dist$ the geodesic distance on the flat torus
$\T_L^n$. For a nonempty closed set $D\subset\T_L^n$, set
\[
  \dist(y,D)\coloneqq\inf_{x\in D}\dist(y,x),
  \qquad
  R(D)\coloneqq\max_{y\in\T_L^n}\dist(y,D).
\]
There exists $\varepsilon_0>0$, depending only on
$d$, $L$, and $\|V\|_{C^2}$, such that every
$\varepsilon$-outward barrier certificate with
$0\le\varepsilon\le\varepsilon_0$ satisfies
\[
  R(D)\le2d\varepsilon.
\]
In particular, every exact outward barrier certificate is the whole
torus.
\end{theorem}

\begin{proof}[Proof of \Cref{thm:no barrier}]
If $D=\T_L^n$, the conclusion is immediate. We therefore assume that
$D$ is a proper subset of $\T_L^n$. Let
  $C_2\coloneqq 1/2\|D^2V\|_\infty$.
Choose $r_0>0$ such that
\[
  r_0<\min\left\{d,\frac L2\right\},
  \qquad
  C_2r_0\le\frac1{2d},
\]
with the second restriction omitted if $C_2=0$.  Choose $0<\e_0<r_0/(2d)$.
Take $y\notin D$ with
  $r\coloneqq \dist(y,D)<r_0$.
Let $x\in D$ be a nearest point and write, in compatible flat
coordinates,
\[
  y=x+r\nu,
  \qquad |\nu|=1.
\]
Then $\nu$ is a proximal outward normal. Set
\[
  s\coloneqq
  \nu^{\mathsf T}S_V(x)\nu
  =
  \nu^{\mathsf T}DV(x)\nu.
\]
The barrier condition gives
\[
  V(x)\cdot \nu
  \le\e-(1+ds)_+.
\]
Taylor's theorem yields
\begin{align*}
  V(y)\cdot \nu
  \le V(x)\cdot \nu+rs+C_2r^2\le\e-(1+ds)_++rs+C_2r^2.
\end{align*}
The elementary cutoff inequality
\begin{equation*}
  -(1+ds)_++rs\le-\frac rd,
  \qquad 0<r<d,
\end{equation*}
holds for every $s\in\R$.  Indeed, if $s\le-1/d$, the left side is $rs\le-r/d$; if $s>-1/d$, it equals $-1-(d-r)s$, whose supremum is $-r/d$.  Therefore,
\begin{equation}\label{eq:drift-app}
  V(y)\cdot \nu
  \le\e-\frac r{2d}.
\end{equation}

Let $\Phi_t$ be the flow generated by $V$, and set 
$\rho(t)\coloneqq\dist(\Phi_t(y),D)$. 
Whenever $0<\rho(t)<r_0$,
\Cref{lem:distance-Dini-background} and \eqref{eq:drift-app},
applied at the current point $\Phi_t(y)$, give
\[
  \Dplus\rho(t)
  \le\varepsilon-\frac{\rho(t)}{2d}.
\]
Let $I$ be a connected component of
$\{t\ge0:\rho(t)>0\}$. If $0\in I$, write $I=[0,\beta)$;
otherwise write $I=(\alpha,\beta)$, where $\alpha>0$ and continuity
gives $\rho(\alpha)=0$. On any portion of $I$ on which
$\rho<r_0$, scalar comparison for the upper Dini inequality gives
\[
  \rho(t)
  \le
  2d\varepsilon+
  \bigl(\rho(0)-2d\varepsilon\bigr)e^{-t/(2d)},
  \qquad 0\in I,
\]
and
\[
  \rho(t)
  \le
  2d\varepsilon
  \bigl(1-e^{-(t-\alpha)/(2d)}\bigr)
  \le2d\varepsilon,
  \qquad 0\notin I.
\]
In the first case the right hand side is a convex combination of
$\rho(0)$ and $2d\varepsilon$; in the second it is at most
$2d\varepsilon$. Since both $\rho(0)$ and $2d\varepsilon$ are
strictly smaller than $r_0$, a first exit argument extends the
corresponding estimate throughout $I$. Using continuity at the finite
endpoints of the components, and observing that $\rho=0$ outside their
union, we obtain
\begin{equation}\label{eq:rho-app}
  \rho(t)
  \le
  2d\varepsilon+
  \bigl(\rho(0)-2d\varepsilon\bigr)_+e^{-t/(2d)},
  \qquad t\ge0.
\end{equation}
The same estimate also holds for trajectories starting at a point
$y\in D$. In that case no connected component of
$\{t\ge0:\rho(t)>0\}$ contains $0$, and the second comparison above
gives
  $\rho(t)\le 2d\varepsilon$, 
  for all $t\ge0$.
Thus \eqref{eq:rho-app} is valid for every initial point satisfying
$\rho(0)<r_0$.
Suppose that $R(D)>2d\varepsilon$. Choose
\[
  2d\varepsilon<r<\min\{R(D),r_0\}
\]
and write
  $D_s\coloneqq\{y\in\T_L^n:\dist(y,D)\le s\}$.
Equation \eqref{eq:rho-app} gives
\[
  \Phi_t(D_r)\subset D_{r_t},
  \qquad
  r_t=2d\varepsilon+
      (r-2d\varepsilon)e^{-t/(2d)}.
\]
Since $V$ is divergence free, its flow preserves Lebesgue measure.
Therefore
\[
  |D_r|
  =
  |\Phi_t(D_r)|
  \le
  |D_{r_t}|.
\]
As $r_t\downarrow2d\varepsilon$, continuity from above of Lebesgue
measure gives
\[
  |D_r|\le|D_{2d\varepsilon}|.
\]
But $0\le2d\varepsilon<r<R(D)$, and the final observation in
Appendix~\ref{app:proximal-background} shows that
\[
  \{y:2d\varepsilon<\dist(y,D)<r\}
\]
is a nonempty open set and hence has positive measure. Thus
  $|D_{2d\varepsilon}|<|D_r|$,
a contradiction. Therefore $R(D)\le2d\varepsilon$.
\end{proof}

\end{document}
